\section{Preliminaries and unconditional starting theorems}
\label{sec:preliminaries}

All finite-set averages and Fourier transforms use probability normalization.
Every affine space in the paper is a finite nonempty affine space over $\F_2$.
Constants described as fixed are independent of the input instance; their
dependence on the requested approximation factor is allowed.

\subsection{Affine alphabets and odd characters}

For an affine space $A$, let $D(A)$ be the vector space of affine functions
$A\to\F_2$, including constants. If $A$ has dimension $r$, then $D(A)$ has
dimension $r+1$. Write $1\in D(A)$ for the constant-one function. A frequency
is an element $\alpha$ of the linear dual $D(A)^*$, with character
$\chi_\alpha(b)=(-1)^{\alpha(b)}$. Call $\alpha$ \emph{odd} if $\alpha(1)=1$.

\begin{lemma}[Odd frequencies are labels]\label{lem:odd-labels}
The map $a\mapsto\operatorname{ev}_a$, where
$\operatorname{ev}_a(b)=b(a)$, is a bijection from $A$ onto the odd frequencies
of $D(A)$. If $\pi:A\to A'$ is affine and
$\pi^*:D(A')\to D(A)$ is pullback, then its dual sends
$\operatorname{ev}_a$ to $\operatorname{ev}_{\pi(a)}$.
\end{lemma}
\begin{proof}
Choose affine coordinates $A=\F_2^r$. Every affine function is
$b(x)=b_0+\sum_{j=1}^r b_jx_j$. A functional taking the constant function to
one is uniquely $b\mapsto b_0+\sum_j b_ja_j$, for some $a\in\F_2^r$.
Finally $(\pi^*)^*(\operatorname{ev}_a)(b)=b(\pi(a))$.
\end{proof}

For a real function $v$ on a finite binary vector space $W$, write
\[
 \widehat v(\alpha)=\E_{w\in W}v(w)(-1)^{\alpha(w)},\qquad
 v(w)=\sum_{\alpha\in W^*}\widehat v(\alpha)(-1)^{\alpha(w)}.
\]
Character orthogonality gives Parseval's identity
$\sum_\alpha\abs{\widehat v(\alpha)}^2=\E v^2$.
If $T:W'\to W$ is linear, each character of $v$ pulls back to its image
under $T^*$; coefficients with the same image add. Injectivity and
surjectivity of $T$ are unnecessary. We will use the identity
\begin{equation}\label{eq:fourth-moment}
 \sum_\alpha\abs{\widehat v(\alpha)}^4
 =\E_{a,b,c\in W}v(a)v(b)v(c)v(a+b+c).
\end{equation}
Indeed, expansion of the four factors and averaging over the three independent
points force all four frequencies to agree.

\subsection{Gaussian identities}

We recall the Fourier--Hermite identities used in Gaussian
noise-stability analysis~\cite[Sections~4.1 and~5]{MOO10}, with proofs of the needed forms.
Let $\gamma_d$ be standard Gaussian probability measure on $\R^d$. The normalized
Hermite polynomials $H_j$ are specified by
\[
 e^{tz-t^2/2}=\sum_{j\ge0}H_j(z)\frac{t^j}{\sqrt{j!}}.
\]
Gaussian integration of the product of two generating functions proves their
orthonormality. Their products
$\psi_I(c)=\prod_x H_{I_x}(c_x)$ form an orthonormal basis of
$L^2(\gamma_d)$. For completeness, if a function is orthogonal to all these
polynomials, its Gaussian-weighted exponential transform has every Taylor
coefficient zero. Cauchy--Schwarz makes that transform entire in each variable;
its restriction to imaginary arguments is the Fourier transform of an
integrable function, which is therefore zero. Thus the original function is
zero almost everywhere.

Differentiating the generating function gives
$H_j'=\sqrt j H_{j-1}$. Consequently, for a function with square-integrable
weak gradient,
\begin{equation}\label{eq:hermite-energy}
 \int\norm{\nabla F}^2\dd\gamma_d
 =\sum_I |I|\abs{\widehat F(I)}^2,
 \qquad |I|=\sum_x I_x.
\end{equation}
One first proves this for polynomials, and then passes to the Hermite partial
sums and weak derivatives. In the smoothed functions used below the weighted
series converges directly. In particular, mean-zero $F$ satisfies the Gaussian
Poincar\'e inequality $\int F^2\dd\gamma_d\le\int\norm{\nabla F}^2\dd\gamma_d$.

For $0<\rho<1$, the Gaussian noise operator
\[
 T_\rho Q(c)=\E_g Q(\rho c+\sqrt{1-\rho^2}\,g)
\]
multiplies $\psi_I$ by $\rho^{|I|}$, as follows by averaging its generating
function. If $Q$ is bounded and $\sigma>0$, ordinary Gaussian convolution is
smooth and
\begin{equation}\label{eq:gaussian-derivative}
 \partial_{c_x}\E_g Q(c+\sigma g)
 =\frac1\sigma\E_g g_xQ(c+\sigma g).
\end{equation}
Differentiation of the Gaussian density is justified by its integrable
derivatives; an additional independent average may be included in $Q$.

We also use the elementary sign identity underlying random-hyperplane
rounding~\cite[Lemma~3.2]{GW95}: if $u,v$ are unit vectors and $g$ is
standard isotropic Gaussian, then
\begin{equation}\label{eq:gaussian-signs}
 \Prob[\sgn\ip gu\ne\sgn\ip gv]
 =\frac{\angle(u,v)}\pi\le\norm{u-v}.
\end{equation}
It follows by rotational symmetry in the two-dimensional plane spanned by
$u,v$; the collinear cases follow by continuity. In particular, adding an
independent centered Gaussian of standard deviation $a$ to a standard Gaussian changes
its sign with probability $\arctan(a)/\pi\le a$.

\subsection{The starting gap and repetition}

\begin{theorem}[An unconditional equation gap]\label{thm:max3lin}
For every fixed rational $\zeta>0$, there is a deterministic polynomial-time
reduction from $3$SAT to a nonempty explicit list of equations
$x_i+x_j+x_k=b$ over $\F_2$, with uniform sampling from the list, such that
satisfiable inputs admit an assignment satisfying at least $1-\zeta$ of the
list and unsatisfiable inputs admit none satisfying more than $3/4$.
The list length and the full bit complexity are polynomial for fixed $\zeta$.
\end{theorem}
\begin{proof}
H\aa stad's Theorem~5.4, including the verifier-to-equations construction in
its proof, gives a deterministic polynomial construction of rationally weighted
three-variable equations with YES value at least $1-\delta$ and NO value at
most $(1+\delta)/2$, for every fixed sufficiently small dyadic
$\delta>0$~\cite{Hastad2001}. Here the weights and their polynomial-size
support are explicit.

We record why a uniform list costs only polynomial work. For a probability
vector $(w_1,\ldots,w_r)$ of rational weights and an integer $Q$, take
$\lfloor Qw_j\rfloor$ copies of equation $j$ and distribute the remaining
copies among the largest fractional remainders. The resulting probabilities
differ from the original ones in total variation by at most $r/(2Q)$.
All floors, comparisons and sorting require polynomial bit complexity.
Choose a fixed $\eps>0$ and $\delta>0$ with
$\delta+\eps<\zeta$ and $(1+\delta)/2+\eps<3/4$, then
$Q\ge r/(2\eps)$ with $Q=O_\eps(r)$. Satisfaction probabilities change by at
most $\eps$ for every assignment simultaneously, giving the asserted list.
The underlying construction can be taken nonempty; any exceptional bounded
input cases are decidable directly and replaced by fixed YES or NO lists.
\end{proof}

A finite two-player game consists of a joint distribution on questions
$(U,V)$, finite answer sets, and an acceptance predicate. Its value is the
maximum acceptance probability of separate deterministic response functions.
Independent shared or private randomness does not increase this maximum:
fixing the random seed realizes at least its average success. The ordinary
$r$-fold repetition samples $r$ independent question pairs and accepts only
if all coordinates accept.

\begin{theorem}[Uniform parallel repetition]\label{thm:parallel-repetition}
There is an effectively specified constant $c<1$ such that every finite game
with at most eight first-player answers, at most two second-player answers,
and value at most $11/12$ has $r$-fold repeated value at most $c^r$ for all
integers $r\ge0$. The same $c$ works for all question distributions and all
question-set sizes.
\end{theorem}
\begin{proof}
This is an instance of ordinary parallel repetition~\cite{Raz1998}.
More explicitly, Theorem~2.5 of Holenstein~\cite{Holenstein2009}, after padding
the answer sets to sizes eight and two, permits
\[
 c=\left(1-\frac1{6000\cdot12^3}\right)^{1/4}<1.
\]
Holenstein's Theorem~2.5 imposes no regularity condition on the joint question
distribution. A rational upper bound strictly between this $c$ and one may
be used in every effective parameter choice below.
\end{proof}
